% Parte 2: equipo. Detalle en docs/team.md.
\docpart{server}{El equipo}{Cuatro roles pensados para la rúbrica, un cronograma por bloques y pocas reglas.}

\section{Roles}
Algoritmo y Datos construyen la defensa (Impacto); Integración la mide y la deja reproducible (Ejecución); Pitch la vuelve entendible en 2 minutos (Presentación).

\begin{card}{algo}{Algoritmo (líder)}
\textbf{Hace:} el paquete \code{guard/}: trozos, embeddings, puntajes, umbrales y saneo.\par
\textbf{Primera hora:} elegir el modelo de embeddings y medir cuánto tarda.\par
\textbf{Listo cuando:} \code{check()} real cumple el contrato, tarda menos de 300 ms y supera al filtro de palabras clave en validación.
\end{card}

\begin{card}{datos}{Datos}
\textbf{Hace:} inyecciones escritas como salidas de herramientas, salidas legítimas y \emph{hard negatives} (textos legítimos que dan órdenes, como “please send the report by Friday”), separados por fuente.\par
\textbf{Primera hora:} fuentes con su licencia en \code{CITATIONS.md} y las primeras 20 filas.\par
\textbf{Listo cuando:} hay al menos 150 inyecciones y 150 legítimos, \textbf{sin ningún texto de AgentDojo}.
\end{card}

\begin{card}{server}{Integración y benchmarks}
\textbf{Hace:} el repo ejecutable, AgentDojo con la clave de OpenAI, el adaptador del guardián, las líneas base, \code{results/}, el README y la demo para la ficha.\par
\textbf{Primera hora:} instalar AgentDojo y correr una tarea sin defensa.\par
\textbf{Listo cuando:} un solo comando del README reproduce la tabla de la ficha.
\end{card}

\begin{card}{pitch}{Pitch}
\textbf{Hace:} registro, ficha técnica, resumen, guion, video de la demo, citas y entrega. Va a todos los checkpoints.\par
\textbf{Primera hora:} \textbf{registrar al equipo antes de las 11:00}, pedir la plantilla oficial de la ficha y bocetarla.\par
\textbf{Listo cuando:} la ficha tiene los números de \code{results/} y todo se entregó a las 16:00.
\end{card}

Si son cinco, la quinta persona hace un proxy MCP que pasa las respuestas de cualquier servidor por \code{check()}: vuelve la idea adoptable, pero nunca se arriesga la evaluación por esto.

\clearpage
\section{Cronograma (sábado 14 de noviembre)}
\begin{center}
\small
\begin{tabularx}{\linewidth}{>{\raggedright}p{24mm}>{\raggedright}X>{\raggedright\arraybackslash}X}
\toprule
\textbf{Bloque} & \textbf{Qué se hace} & \textbf{Checkpoint}\\\midrule
1 · 9:00–10:00 & modelo de embeddings; fuentes de datos; AgentDojo instalado; \textbf{registro} y plantilla & —\\
2 · 10:00–10:45 & \code{check()} v0; formato de datos; adaptador con el mock; guion preliminar & \textbf{CP1:} AgentDojo corre una tarea con el mock\\
3 · 10:45–13:00 & \code{check()} v1; dataset completo; líneas base completas; ficha v1 & \textbf{CP2:} guardián real en AgentDojo; líneas base en \code{results/}\\
4 · 13:00–15:00 & umbrales en validación y congelar; corridas finales; demo y capturas; resultados en la ficha & \textbf{CP3:} números finales en la ficha\\
5 · 15:00–16:00 & README reproducible, citas completas, límites y uso dual & \textbf{entrega a las 16:00}\\
6 · 16:00–17:00 & dos ensayos con la ficha en pantalla & pitch a las 17:10\\\bottomrule
\end{tabularx}
\end{center}
El reparto por rol de cada bloque está en \Docs{team.md}.

\section{Reglas}
\begin{enumerate}
  \item En cada checkpoint, cada persona muestra algo \textbf{funcionando}.
  \item Nadie se queda atascado más de 30 minutos sin preguntar.
  \item Primero la versión fea de punta a punta, después se mejora.
  \item Los contratos no cambian sin avisar en un checkpoint.
  \item Todo lo que se usa (código, datos, texto generado con IA) se cita en \code{CITATIONS.md} en el momento.
  \item \textbf{De la hackathon:} nada contra sistemas reales; los ataques son datos, nunca se ejecutan ni se publican como herramienta; solo datos públicos o sintéticos.
\end{enumerate}
